\chapter*{Abstract}
\addcontentsline{toc}{chapter}{Abstract}

Tensor programs need a source boundary that users can write directly and
compilers can inspect. A common pattern still crosses several boundaries:
multidimensional indexing, recurrent state updates, and sensitivity queries are
usually expressed through separate interfaces such as \texttt{einsum}-like
functions, host-language loops or scan combinators, and \texttt{grad} or
\texttt{backward} APIs. This fragmentation makes programs executable, but it
hides the source facts that connect the mathematical definition to compiler
analysis.

This thesis presents \einlang{}, a tensor language and compiler that integrates
multidimensional indexed definitions, declarative recurrences, and local
derivative requests in one binding environment. In \einlang{}, a recurrence is a
named binding over an index space that declares its dependencies; later
expressions can read points of that recurrence, and derivative requests can
target either final results or intermediate bindings. The contribution is an
executable abstraction boundary: indexed clauses define axes and covered
domains, recurrent bindings define dependency graphs, and derivative requests
denote shaped derivative values.

The thesis also documents the non-trivial implementation pieces that make this
boundary executable: the grammar and transformer front end, DefId-based name
resolution, module loading and standard-library discovery, the IR and pass
pipeline, range and shape analyses, type inference and monomorphization,
Einstein lowering, recurrence ordering and storage metadata, compile-time and
runtime autodiff machinery, NumPy execution, the experimental IREE backend,
diagnostics, testing, and example validation. The result is a source form for
tensor programs that stays close to the mathematical description while
preserving structure for compiler and runtime passes.
